\documentclass[../DoAn.tex]{subfiles}
\begin{document}


\chapter{GIỚI THIỆU}
% Lưu ý: \textbf{Trước khi viết ĐATN, sinh viên cần đọc kỹ hướng dẫn và quy định chi tiết} về cách viết ĐATN trong Phụ lục A. Sinh viên tuân theo mẫu tài liệu này để viết báo cáo đồ án tốt nghiệp, vì tài liệu này đã được căn chỉnh, chỉnh sửa theo đúng chuẩn báo cáo kỹ thuật đồ án tốt nghiệp (ISO 7144:1986). Sinh viên viết trực tiếp vào file này, chỉ chỉnh sửa nội dung, và không viết trên file mới.

% \textbf{Khi đóng quyển ĐATN}, sinh viên cần lưu ý tuân thủ hướng dẫn ở phụ lục A.9

% \textbf{SV cần đặc biệt lưu ý cách hành văn}. Mỗi đoạn văn không được quá dài và cần có ý tứ rõ ràng, bao gồm duy nhất một ý chính và các ý phân tích bổ trợ để làm rõ hơn ý chính. Các câu văn trong đoạn phải đầy đủ chủ ngữ vị ngữ, cùng hướng đến chủ đề chung. Câu sau phải liên kết với câu trước, đoạn sau liên kết với đoạn trước. Trong văn phong khoa học, sinh viên không được dùng từ trong văn nói, không dùng các từ phóng đại, thái quá, các từ thiếu khách quan, thiên về cảm xúc, về quan điểm cá nhân như “tuyệt vời”, “cực hay”, “cực kỳ hữu ích”, v.v. Các câu văn cần được tối ưu hóa, đảm bảo rất khó để thể thêm hoặc bớt đi được dù chỉ một từ. Cách diễn đạt cần ngắn gọn, súc tích, không dài dòng.

% Mẫu ĐATN này được thiết kế phù hợp nhất với đa số các đề tài xây dựng phần mềm ứng dụng. Với các dạng đề tài khác (giải pháp, nghiên cứu, phần mềm đặc thù, v.v.), sinh viên dựa trên cấu trúc và hướng dẫn của báo cáo này để đề xuất và trao đổi với giáo viên hướng dẫn để thiết kế khung báo cáo đồ án cho phù hợp. Sinh viên lưu ý \textbf{trong mọi trường hợp, SV luôn phải sử dụng định dạng báo cáo này, và phải đọc kỹ toàn bộ các hướng dẫn từ đầu tới cuối.} Các hướng dẫn không chỉ áp dụng riêng cho đề tài ứng dụng, mà còn phù hợp với các dạng đề tài khác. Ngoài ra, trong mẫu ĐATN này đã được tích hợp một số hướng dẫn dành riêng cho đề tài nghiên cứu.

% Chương 1 có độ dài từ 3 đến 6 trang với các nội dung sau đây

\section{Đặt vấn đề}
\label{section:1.1}

Trong bối cảnh chuyển đổi số, thư điện tử đã trở thành một trong những phương tiện giao tiếp quan trọng trong học tập và công việc. Hằng ngày, người dùng phải tiếp nhận nhiều loại email khác nhau như thông báo, yêu cầu công việc, lời mời họp, báo cáo tiến độ hay các thư trao đổi với đối tác và khách hàng. Sự gia tăng về số lượng email khiến việc quản lý, theo dõi và xử lý thông tin ngày càng trở nên khó khăn.

Bên cạnh chức năng trao đổi thông tin, email còn chứa nhiều dữ liệu liên quan đến lịch làm việc, lịch họp và các nhiệm vụ cần thực hiện. Người dùng thường phải dành nhiều thời gian để đọc, phân loại, tìm kiếm nội dung quan trọng và theo dõi các yêu cầu cần phản hồi. Việc bỏ sót hoặc chậm xử lý các email quan trọng có thể ảnh hưởng trực tiếp đến hiệu quả công việc, tiến độ dự án và khả năng phối hợp giữa các bên liên quan.

Do đó, nhu cầu hỗ trợ quản lý và xử lý email một cách hiệu quả đang trở nên ngày càng cấp thiết. Nếu bài toán này được giải quyết tốt, người dùng có thể giảm đáng kể thời gian dành cho các thao tác thủ công, nâng cao năng suất làm việc và cải thiện khả năng quản lý thời gian. Ngoài ra, kết quả nghiên cứu còn có thể được mở rộng và ứng dụng trong nhiều lĩnh vực khác như quản lý công việc, hỗ trợ khách hàng, giáo dục và các hệ thống trợ lý số thông minh.


\section{Mục tiêu và phạm vi đề tài}
\label{section:1.2}

Hiện nay, nhiều hệ thống thư điện tử như Gmail, Outlook hay Yahoo Mail đã cung cấp các tính năng hỗ trợ người dùng quản lý hộp thư như phân loại thư theo danh mục, tìm kiếm nội dung, gắn nhãn hoặc lọc thư rác. Bên cạnh đó, sự phát triển của trí tuệ nhân tạo đã thúc đẩy sự xuất hiện của các công cụ hỗ trợ soạn thảo và trả lời email tự động, giúp người dùng tiết kiệm thời gian trong quá trình xử lý công việc.

Mặc dù vậy, phần lớn các hệ thống hiện nay vẫn yêu cầu người dùng chủ động đọc và xử lý từng email, đồng thời phải tự theo dõi các thông tin liên quan đến lịch làm việc, lịch họp và các yêu cầu cần phản hồi. Các chức năng hỗ trợ thường được triển khai rời rạc, chưa tạo thành một môi trường tương tác thống nhất giúp người dùng quản lý và khai thác thông tin từ email một cách thuận tiện. Đối với những người thường xuyên làm việc với số lượng lớn email, việc tổng hợp thông tin, theo dõi các email quan trọng và thực hiện các thao tác xử lý vẫn tiêu tốn nhiều thời gian.

Xuất phát từ những hạn chế trên, đề tài hướng đến xây dựng một hệ thống trợ lý thông minh hỗ trợ quản lý và khai thác thông tin từ email. Hệ thống tập trung vào các chức năng chính như: thu thập và phân loại email, hỗ trợ tra cứu và tóm tắt thông tin qua giao diện hội thoại, hỗ trợ quản lý lịch làm việc thông qua dữ liệu email, đồng thời hỗ trợ người dùng soạn thảo và phản hồi email theo yêu cầu. Mục tiêu của đề tài là giảm khối lượng công việc thủ công trong quá trình xử lý email, nâng cao hiệu quả quản lý thông tin và cải thiện trải nghiệm làm việc của người dùng.

Phạm vi của đề tài tập trung vào việc xử lý email cá nhân thông qua dịch vụ Gmail và xây dựng giao diện tương tác dưới dạng trợ lý hội thoại. Đề tài không đi sâu vào các bài toán bảo mật nâng cao, tối ưu hạ tầng triển khai quy mô lớn hoặc xây dựng mô hình ngôn ngữ mới, mà tập trung vào việc ứng dụng các công nghệ hiện có để giải quyết bài toán hỗ trợ quản lý email trong thực tế.


\section{Định hướng giải pháp}
\label{section:1.3}

Để giải quyết các vấn đề đã trình bày ở phần \ref{section:1.2}, đề tài định hướng xây dựng một hệ thống trợ lý email thông minh dựa trên sự kết hợp giữa công nghệ xử lý ngôn ngữ tự nhiên, mô hình ngôn ngữ lớn và các dịch vụ API của Gmail và Google Calendar. Việc lựa chọn các công nghệ này nhằm tận dụng khả năng hiểu ngữ nghĩa, tổng hợp thông tin và tương tác hội thoại của các mô hình AI hiện đại, đồng thời khai thác trực tiếp dữ liệu email và lịch làm việc của người dùng thông qua các nền tảng phổ biến.

Trên cơ sở đó, đề tài đề xuất xây dựng một hệ thống trợ lý hội thoại cho phép người dùng tương tác bằng ngôn ngữ tự nhiên để tóm tắt và khai thác thông tin từ email. Hệ thống có khả năng thu thập và quản lý email từ Gmail, hỗ trợ tìm kiếm thông tin liên quan, phát hiện các sự kiện hoặc cuộc họp được đề cập trong email, hỗ trợ quản lý lịch làm việc, đồng thời hỗ trợ soạn thảo và phản hồi email theo yêu cầu của người dùng. Các tác vụ được thực hiện thông qua một giao diện hội thoại thống nhất nhằm giảm số lượng thao tác thủ công trong quá trình xử lý email.

Đóng góp chính của đề tài là xây dựng và tích hợp thành công một hệ thống trợ lý email thông minh kết hợp giữa quản lý email, hỗ trợ lịch làm việc và giao diện hội thoại AI trong cùng một nền tảng. Kết quả đạt được là một hệ thống có khả năng hỗ trợ người dùng tiếp cận và xử lý thông tin email hiệu quả hơn, giảm thời gian tìm kiếm và tổng hợp thông tin, đồng thời nâng cao trải nghiệm quản lý công việc thông qua tương tác bằng ngôn ngữ tự nhiên.


\section{Bố cục đồ án}
\label{section:1.4}


% Phần còn lại của báo cáo đồ án tốt nghiệp này được tổ chức như sau. 

% Chương 2 trình bày về v.v. 

% Trong Chương 3, em/tôi giới thiệu về v.v.

% \textbf{Chú ý:} Sinh viên cần viết mô tả thành đoạn văn đầy đủ về nội dung chương. Tuyệt đối không viết ý hay gạch đầu dòng. Chương 1 không cần mô tả trong phần này. 

% Ví dụ tham khảo mô tả chương trong phần bố cục đồ án tốt nghiệp: Chương *** trình bày đóng góp chính của đồ án, đó là một nền tảng ABC cho phép khai phá và tích hợp nhiều nguồn dữ liệu, trong đó mỗi nguồn dữ liệu lại có định dạng đặc thù riêng. Nền tảng ABC được phát triển dựa trên khái niệm DEF, là các module ngữ nghĩa trợ giúp người dùng tìm kiếm, tích hợp và hiển thị trực quan dữ liệu theo mô hình cộng tác và mô hình phân tán.

% \textbf{Chú ý:} Trong phần nội dung chính, mỗi chương của đồ án nên có phần Tổng quan và Kết chương. Hai phần này đều có định dạng văn bản “Normal”, sinh viên không cần tạo định dạng riêng, ví dụ như không in đậm/in nghiêng, không đóng khung, v.v. 

% Trong phần Tổng quan của chương N, sinh viên nên có sự liên kết với chương N-1 rồi trình bày sơ qua lý do có mặt của chương N và sự cần thiết của chương này trong đồ án. Sau đó giới thiệu những vấn đề sẽ trình bày trong chương này là gì, trong các đề mục lớn nào.

% Ví dụ về phần Tổng quan: Chương 3 đã thảo luận về nguồn gốc ra đời, cơ sở lý thuyết và các nhiệm vụ chính của bài toán tích hợp dữ liệu. Chương 4 này sẽ trình bày chi tiết các công cụ tích hợp dữ liệu theo hướng tiếp cận “mashup”. Với mục đích và phạm vi của đề tài, sáu nhóm công cụ tích hợp dữ liệu chính được trình bày bao gồm: (i) nhóm công cụ ABC trong phần 4.1, (ii) nhóm công cụ DEF trong phần 4.2, nhóm công cụ GHK trong phần 4.3, v.v.

% Trong phần Kết chương, sinh viên đưa ra một số kết luận quan trọng của chương. Những vấn đề mở ra trong Tổng quan cần được tóm tắt lại nội dung và cách giải quyết/thực hiện như thế nào. Sinh viên lưu ý không viết Kết chương giống hệt Tổng quan. Sau khi đọc phần Kết chương, người đọc sẽ nắm được sơ bộ nội dung và giải pháp cho các vấn đề đã trình bày trong chương. Trong Kết chương, Sinh viên nên có thêm câu liên kết tới chương tiếp theo.

% Ví dụ về phần Kết chương: Chương này đã phân tích chi tiết sáu nhóm công cụ tích hợp dữ liệu. Nhóm công cụ ABC và DEF thích hợp với những bài toán tích hợp dữ liệu phạm vi nhỏ. Trong khi đó, nhóm công cụ GHK lại chứng tỏ thế mạnh của mình với những bài toán cần độ chính xác cao, v.v. Từ kết quả nghiên cứu và phân tích về sáu nhóm công cụ tích hợp dữ liệu này, tôi đã thực hiện phát triển phần mềm tự động bóc tách và tích hợp dữ liệu sử dụng nhóm công cụ GHK. Phần này được trình bày trong chương tiếp theo – Chương 5.

\end{document}